\documentclass[fleqn]{article}
\oddsidemargin 0.0in
\textwidth 6.0in
\thispagestyle{empty}
\usepackage{import}
\usepackage{amsmath}
\usepackage{graphicx}
\usepackage{flexisym}
\usepackage{amssymb}
\usepackage{bigints} 
\usepackage[english]{babel}
\usepackage[utf8x]{inputenc}
\usepackage{float}
\usepackage[colorinlistoftodos]{todonotes}

\definecolor{hwColor}{HTML}{AD53BA}

\begin{document}

  \begin{titlepage}

    \newcommand{\HRule}{\rule{\linewidth}{0.5mm}}

    \center


    \textsc{\LARGE Arizona State University}\\[1.5cm]

    \textsc{\LARGE Quantum Physics I }\\[1.5cm]


    \begin{figure}
      \includegraphics[width=\linewidth]{asu.png}
    \end{figure}


    \HRule \\[0.4cm]
    { \huge \bfseries Homework Nine}\\[0.4cm] 
    \HRule \\[1.5cm]

    \textbf{Behnam Amiri}

    \bigbreak

    \textbf{Prof: Richard Kirian}

    \bigbreak


    \textbf{{\large \today}\\[2cm]}

    \vfill 

  \end{titlepage}

  \textbf{2.25} \\ \\
  Check that the bound state of the delta-function well (Equation 2.132) is orthogonal to the 
  scattering states (Equations 2.134 and 2.135).

    \textcolor{hwColor}{
      \\\\
      Equation (2.132) 
      $$\psi(x)=\dfrac{\sqrt{m\alpha}}{\hbar} e^{-m\alpha|x|/\hbar^2}$$
      and equations 2.134 and 2.135 are respectively $\psi(x)=A e^{ikx}+B e^{-ikx}$ and $\psi(x)=F e^{ikx}+G e^{-ikx}$. \\
      \\
      Checking the orthogonality: \\
      \\
      $
        \bigints_{-\infty}^{+\infty} \psi^*_m \psi_n dx=\bigints_{-\infty}^{0} \psi^*_m \psi_n dx+\bigints_{0}^{+\infty} \psi^*_m \psi_n dx \\
        \\
        \\
        =\bigints_{-\infty}^{0} \left(\dfrac{\sqrt{m\alpha}}{\hbar} e^{m\alpha x/\hbar^2}\right)^* \left(A e^{ikx}+B e^{-ikx}\right) dx
        +\bigints_{0}^{+\infty} \left(\dfrac{\sqrt{m\alpha}}{\hbar} e^{-m\alpha x/\hbar^2}\right)^* \left(F e^{ikx}+G e^{-ikx}\right) dx \\
        \\
        \\
        =\bigints_{-\infty}^{0} \left(\dfrac{\sqrt{m\alpha}}{\hbar} e^{m\alpha x/\hbar^2}\right) \left(A e^{ikx}+B e^{-ikx}\right) dx
        +\bigints_{0}^{+\infty} \left(\dfrac{\sqrt{m\alpha}}{\hbar} e^{-m\alpha x/\hbar^2}\right) \left(F e^{ikx}+G e^{-ikx}\right) dx \\
        \\
        \\
        =\dfrac{\sqrt{m\alpha}}{\hbar} \left[\bigints_{-\infty}^{0} \left(e^{m\alpha x/\hbar^2}\right) \left(A e^{ikx}+B e^{-ikx}\right) dx
        +\bigints_{0}^{+\infty} \left(e^{-m\alpha x/\hbar^2}\right) \left(F e^{ikx}+G e^{-ikx}\right) dx \right] \\
        \\
        \\
        =\dfrac{\sqrt{m\alpha}}{\hbar} (
          \dfrac{A}{\dfrac{m \alpha}{\hbar^2}+ik}e^{(m\alpha/\hbar^2+ik)x}\Big|_{-\infty}^{0}
          +\dfrac{B}{\dfrac{m \alpha}{\hbar^2}-ik}e^{(m\alpha/\hbar^2-ik)x}\Big|_{-\infty}^{0} \\
          +\dfrac{F}{-\dfrac{m \alpha}{\hbar^2}+ik}e^{(-m\alpha/\hbar^2+ik)x}\Big|_{-\infty}^{0}
          +\dfrac{G}{-\dfrac{m \alpha}{\hbar^2}-ik}e^{(-m\alpha/\hbar^2-ik)x}\Big|_{-\infty}^{0}
        ) \\
        \\
        \\
        =\dfrac{\sqrt{m\alpha}}{\hbar} \left[\dfrac{\dfrac{m \alpha}{\hbar^2}(A+B+F+G)+(-A+B+F-G)ik}{k^2+\left(\dfrac{m \alpha}{h^2}\right)^2}\right] \\
        \\
        \\
      $
      From page 66 of the textbook (Griffiths), the continuity of $\psi(x)$ at $x=0$ requires that: $F+G=A+B$, hence: \\
      \\
      \\
      $
        \Rightarrow \bigints_{-\infty}^{+\infty} \psi^*_m \psi_n dx=\dfrac{\sqrt{m\alpha}}{\hbar} \left[\dfrac{\dfrac{2m \alpha}{\hbar^2}(A+B)-\dfrac{2m \alpha}{\hbar^2}(A+B)}{k^2+\left(\dfrac{m \alpha}{h^2}\right)^2}\right] \\
        \\
        \\
        \\
        \therefore ~~~ \bigints_{-\infty}^{+\infty} \psi^*_m \psi_n dx=0 ~ \surd \\ \\
      $
      Since we got a zero, we can tell that the delta-function well is orthogonal to the scattering states. \\
    }

  \rule{15cm}{1pt}

  \textbf{2.27}
  Consider the \emph{double delta-function} potential \\
  $$V(x)=-\alpha \left[\delta(x+a)+ \delta(x-a)\right]$$
  where $\alpha$ and $\beta$ are positive constants. \\ \\
  \begin{itemize}
    \item Sketch this potential.
    

    \item How many bound states does it possess? Find the allowed energies, for 
    $\alpha =\dfrac{\hbar^2}{ma}$ and for $\alpha=\dfrac{\hbar^2}{4ma}$, and sketch 
    the wave functions.

    \item What are the bound state energies in the limiting cases $(i) ~ a \rightarrow  0$
    and $(ii) ~ a \rightarrow \infty$ (holding $\alpha$ fixed?) Explain why your answers are reasonable, by comparison with the single
    delta-function well.    
  \end{itemize}


  \rule{15cm}{1pt}

  \textbf{2.34}
  Consider the "step" potential: 
  $$V(x)=\begin{cases}
    0, ~~~ x\leqslant 0 \\
    V_0, ~~~ x > 0
  \end{cases}$$
\begin{itemize}
  \item Calculate the reflection coefficient, for the case $E>V_0$, and comment on the answer.

  \item Calculate the reflection coefficient for the case $E>V_0$.

  \item For a potential (such as this one) that does not go back to zero to the right of the barrier,
  the transmission coefficient is not simply $|F|^2/|A|^2$ (with A the incident amplitude and F
  the transmitted amplitude), because the transmitted wave travels at a different speed. Show
  that
  $$T=\sqrt{\dfrac{E-V_0}{E}}\dfrac{|F|^2}{|A|^2}$$
  for $E>V_0$. \emph{Hint:} You can figure it out using Equation 2.99, or—more elegantly, but less
  informatively—from the probability current (Problem 2.18). What is $T$, $E<V_0?$

  \item For $E>V_0$, calculate the transmission coefficient for the step potential, and check that $T+R=1$.
  
\end{itemize}

\rule{15cm}{1pt}

\textbf{2.53 (part a only)}
\textbf{The Scattering Matrix.} The theory of scattering generalizes in a pretty obvious way to
arbitrary localized potentials (Figure 2.21). To the left (Region I), $V(x)=0$, so 
$$\psi(x)=A e^{ikx}+B e^{-ikx}$$
where $k\equiv \dfrac{\sqrt{2mE}}{\hbar}$
\\
\\
To the right (Region III), V (x) is again zero, so 
$$\psi(x)=F e^{ikx}+G e^{-ikx}$$ 
In between (Region II), of course, I can’t tell you what $\psi$ is until you specify the potential, but
because the Schr$\ddot{o}$dinger equation is a linear, second-order differential equation, the general
solution has got to be of the form 
$$\psi(x)=C f(x)+D g(x)$$,
where f(x) and g(x) are two linearly independent particular solutions. There will be four
boundary conditions (two joining Regions I and II, and two joining Regions II and III). Two of
these can be used to eliminate C and D, and the other two can be “solved” for B and F in terms
of A and G:
$$B=S_{11} A+S_{12} G, ~~~ F=S_{21}A +S_{22} G.$$

The four coefficients Sij , which depend on k (and hence on E), constitute a $2 \times 2$ matrix \textbf{S}, called
the \textbf{scattering matrix} (or \textbf{S-matrix} for short). The S-matrix tells you the outgoing
amplitudes (B and F) in terms of the incoming amplitudes (A and G):
$$\begin{vmatrix}
  B \\
  F
\end{vmatrix}=\begin{vmatrix}
  S_{11} & S_{12} \\
  S_{21} & S_{22}
\end{vmatrix} \begin{vmatrix}
  A \\
  G
\end{vmatrix}.$$
In the typical case of scattering from the left, G = 0, so the reflection and transmission
coefficients are 
$$R_{\ell}=\dfrac{|B|^2}{|A|^2}\Big|=|S_{11}|^2, ~~~~~ T_{\ell}=\dfrac{|F|^2}{|A|^2}\Big|=|S_{21}|^2.$$

For scattering from the right, $A=0$, and 
$$R_r=\dfrac{|F|^2}{|G|^2}\Big|=|S_{22}|^2, ~~~~~ T_r=\dfrac{|B|^2}{|G|^2}\Big|=|S_{12}|^2.$$

\begin{itemize}
  \item Construct the S-matrix for scattering from a delta-function well (Equation 2.117).

\end{itemize}

\end{document}
